\chapter{Master Problem Bank: Advanced Equilibrium Applications of EMF}
\label{chap:qb_equilibrium}

\begin{tcolorbox}[enhanced,colback=subtlegreen,colframe=cengagegreen,arc=2mm,boxrule=1pt,
    title=\textbf{\large Part V Organization: Equilibrium Applications, pH, $K_{sp}$, Latimer \& Frost}]
This part compiles all questions and multi-step analytical problems on solubility products ($K_{sp}$), potentiometric pH determination (SHE, quinhydrone, glass electrode), complex ion formation constants, and Latimer/Frost oxidation state diagrams from all 8 books. Identical and redundant variants from other books are co-located beneath each primary problem with full source citations.
\end{tcolorbox}

\section{Subtopic 5.1: Determination of Solubility Product ($K_{sp}$)}

\begin{problembox}
\textbf{Q.5.1} \hfill \textbf{[Primary Source: Neeraj Kumar Part 1 Ex-II Sec F Q8 / NA Level-1 Q128]}
\label{q:qb-c5-01}

Calculate the solubility product $K_{sp}$ of silver bromide (\ce{AgBr}) at $25^\circ\text{C}$ from the following standard reduction potentials:
\begin{align*}
    \ce{AgBr(s) + e^- -> Ag(s) + Br^-(aq)} \quad \& E^\circ = +0.071\,\text{V} \\
    \ce{Ag+(aq) + e^- -> Ag(s)} \quad \& E^\circ = +0.799\,\text{V}
\end{align*}
\begin{multicols}{4}
\begin{enumerate}[label=(\alph*)]
    \item $5.0 \times 10^{-13}$
    \item $2.5 \times 10^{-12}$
    \item $1.0 \times 10^{-10}$
    \item $7.2 \times 10^{-15}$
\end{enumerate}
\end{multicols}

\tcbline
\textbf{\color{accentamber}Cross-Book Redundant / Identical Variants:}
\begin{itemize}[leftmargin=*,itemsep=2pt]
    \item \textbf{[Variant v1: GRB Practice Problem 50, p. 50]} The EMF of the cell $\ce{Ag} \mid \ce{AgI(s)}, 0.05\,\text{M } \ce{KI} \parallel 0.05\,\text{M } \ce{AgNO3} \mid \ce{Ag}$ is $0.788\,\text{V}$ at $25^\circ\text{C}$. Calculate the solubility product of \ce{AgI}.
    \item \textbf{[Variant v2: Pearson Ex-II Sec E Q18, p. 38]} For the cell $\ce{Pb} \mid \ce{PbSO4(s)}, \ce{SO4^2-}(0.1\,\text{M}) \parallel \ce{Pb^2+}(0.1\,\text{M}) \mid \ce{Pb}$, the EMF is $0.236\,\text{V}$ at $298\,\text{K}$. Calculate $K_{sp}(\ce{PbSO4})$.
    \item \textbf{[Variant v3: Cengage Theory Concept App 5 Q2, p. 43]} How can the solubility product of calomel (\ce{Hg2Cl2}) be determined from standard electrode potentials?
\end{itemize}
\end{problembox}

\section{Subtopic 5.2: Potentiometric $\text{pH}$ Determination (SHE, Quinhydrone \& Glass)}

\begin{problembox}
\textbf{Q.5.2} \hfill \textbf{[Primary Source: GRB  12.28 Solved Ex 48 / NA Level-2 Q36]}
\label{q:qb-c5-02}

A quinhydrone electrode is set up in an unknown acidic buffer solution at $25^\circ\text{C}$ and coupled with a saturated calomel electrode ($E_{\text{SCE}} = +0.242\,\text{V}$). The cell EMF is measured to be $0.180\,\text{V}$ with the quinhydrone electrode as the positive pole. Given $E^\circ(\ce{Q, 2H+/H2Q}) = +0.699\,\text{V}$:
\begin{enumerate}[label=(\alph*)]
    \item Calculate the $\text{pH}$ of the buffer solution.
    \item Why cannot the quinhydrone electrode be used accurately in solutions of $\text{pH} > 8.5$?
\end{enumerate}

\tcbline
\textbf{\color{accentamber}Cross-Book Redundant / Identical Variants:}
\begin{itemize}[leftmargin=*,itemsep=2pt]
    \item \textbf{[Variant v1: Neeraj Kumar Part 1 Ex-II Sec A Q28, p. 14]} The EMF of a cell consisting of a hydrogen electrode dipping in a solution of unknown pH coupled with SCE is $0.450\,\text{V}$. The pH is: (a) $3.52$ (b) $7.61$ (c) $4.25$ (d) $5.80$.
    \item \textbf{[Variant v2: Pearson Ex-II Sec A Q58, p. 21]} In alkaline solution ($\text{pH} > 9$), hydroquinone undergoes: (a) oxidation by dissolved oxygen (b) ionization of phenolic $-OH$ groups ($K_a \approx 10^{-10}$) disrupting the equimolar $[\ce{Q}] = [\ce{H2Q}]$ ratio (c) dimerization (d) precipitation.
    \item \textbf{[Variant v3: Cengage Theory p. 40]} Write the potential equation for a glass electrode and explain how it determines pH across the full range $1 \le \text{pH} \le 12$.
\end{itemize}
\end{problembox}

\section{Subtopic 5.3: Latimer \& Frost Oxidation State Diagrams}

\begin{problembox}
\textbf{Q.5.3} \hfill \textbf{[Primary Source: Atkins Focus 6 / Neeraj Kumar Part 1 Ex-II Sec B Q14]}
\label{q:qb-c5-03}

The Latimer diagram for manganese species in $1\,\text{M}$ acid solution ($\text{pH} = 0$) is:
\begin{equation*}
    \ce{MnO4^-} \xrightarrow{+0.56\,\text{V}} \ce{MnO4^2-} \xrightarrow{+2.27\,\text{V}} \ce{MnO2} \xrightarrow{+0.95\,\text{V}} \ce{Mn^3+} \xrightarrow{+1.51\,\text{V}} \ce{Mn^2+} \xrightarrow{-1.18\,\text{V}} \ce{Mn}
\end{equation*}
\begin{enumerate}[label=(\alph*)]
    \item Calculate the standard potential $E^\circ$ for the direct 5-electron reduction $\ce{MnO4^- + 8H+ + 5e^- -> Mn^2+ + 4H2O}$.
    \item Will manganate ion ($\ce{MnO4^2-}$) undergo spontaneous disproportionation in acid solution? Justify thermodynamically by calculating $E^\circ_{\text{disp}}$.
    \item Will $\ce{Mn^3+}$ undergo spontaneous disproportionation into $\ce{MnO2}$ and $\ce{Mn^2+}$?
\end{enumerate}

\tcbline
\textbf{\color{accentamber}Cross-Book Redundant / Identical Variants:}
\begin{itemize}[leftmargin=*,itemsep=2pt]
    \item \textbf{[Variant v1: Pearson Ex-II Sec A Q68, p. 22]} On a Frost diagram ($N E^\circ$ vs $N$), a species that lies above the line connecting its two neighboring oxidation states is: (a) stable to disproportionation (b) prone to spontaneous disproportionation (c) a reducing agent only (d) in its ground electronic state.
    \item \textbf{[Variant v2: Narendra Avasthi Level-3 Passage 5, p. 39]} Given the standard reduction potentials: $\ce{Fe^3+ + e^- -> Fe^2+}$ ($E^\circ = 0.77\,\text{V}$) and $\ce{Fe^2+ + 2e^- -> Fe}$ ($E^\circ = -0.44\,\text{V}$), calculate $E^\circ$ for $\ce{Fe^3+ + 3e^- -> Fe}$.
    \item \textbf{[Variant v3: Cengage Theory p. 46]} In a Latimer diagram, comproportionation occurs spontaneously between two oxidation states when: $E^\circ_{\text{right}} < E^\circ_{\text{left}}$ for the intermediate state.
\end{itemize}
\end{problembox}
